\section{Introduction}
Evidence in temporal knowledge graphs can change after an answer has been issued. Events from the same observation period may be withdrawn together, leaving a system to decide which answers remain usable. A high prediction score does not reveal whether an answer depends on one such period or survives the removal of several. An explicit tolerance to evidence withdrawal would make that decision inspectable and allow the system to revisit only the affected answers.

\begin{figure*}[!t]
\centering\includegraphics[width=\textwidth]{fig01_interface.pdf}
\caption{\textbf{Evidence organization creates a stability question beyond prediction scores.}
\textbf{a,} A frozen neural anchor and an editable temporal memory provide the distributions used for prediction, calibrated selection and exact deletion certification. Event edits refresh the affected registered queries.
\textbf{b,} A conceptual example with neural probabilities $(0.45,0.35,0.20)$ and evidence weight $0.6$. Both histories have total object masses $(6,4,0)$ and the same score vector $(0.54,0.38,0.08)$, but distribute their support differently across three windows. Distributed rival masses are exactly $4/3$ per window, rounded in the display.
\textbf{c,} Exhaustive deletion of at most the indicated number of windows yields different worst-case winner margins. These values illustrate the construction and are not empirical benchmark results. The schematic and plotting code were prepared with assistance from OpenAI Codex (OpenAI; desktop build 26.915.4065.0 recorded during final preparation).}
\label{fig:interface}
\end{figure*}


Temporal knowledge graph completion has developed increasingly expressive representations of event history. Recent methods model temporal granularity, cross-space interactions and frequency-dependent relations \citep{geng2025granularity,vu2025tcrosse,liu2025terdy}. History extraction, multiple temporal views and state embeddings supply complementary information for predicting missing facts \citep{dao2025hgct,chen2025dualview,ji2025stse}. Newer work extends these representations through Householder transformations, sparse fact aggregation and tensorized temporal relations \citep{xu2026householder,zhu2026splices,lee2026tirano}. These advances improve the predictive model. They also make it valuable to attach an explicit account of how an exposed answer responds to a change in its observable evidence.

Predictive confidence and resistance to evidence deletion address different questions. Calibration estimates how often a prediction is correct, while selective prediction uses a confidence score to determine which answers to expose. Recent uncertainty and rejection methods strengthen this statistical decision layer \citep{liu2025uncertkg,li2025vicinal,szabadvary2025reject,he2025conjunction}. Temporal conformal prediction additionally addresses dependence between graph observations \citep{wang2025nonexchangeable}. An evidence-deletion certificate instead concerns one fixed query: it states that every allowed deletion preserves the answer. Even two queries with the same score vector can have different certificates when their evidence is arranged differently across time (Fig.~\ref{fig:interface}). A useful decision layer should preserve this distinction while making both quantities available at inference time.

We introduce \StableKG{}, an editable evidence interface for neural temporal knowledge graph completion. A frozen neural predictor supplies an anchor distribution, and a temporal memory supplies a second distribution over candidate objects. The memory is renormalized after evidence is deleted. Although this normalization couples all object scores, the sign of each post-deletion winner--competitor margin reduces to an additive expression. Sorting its positive window contributions identifies an exact worst-case deletion without enumerating subsets. A further reduction checks only evidence-supported objects and one additional neural competitor, while a reverse index restricts maintenance to queries affected by an event edit.

The resulting system provides three concrete capabilities: an exact certificate for coherent temporal-window deletion, a selective policy that combines the certificate with calibrated correctness, and local maintenance of predictions and certificates. We evaluate the interface with two temporal neural backbones and three training seeds on ICEWS14 and ICEWS05-15, then examine maintenance on a 2.3-million-entry GDELT history. The experiments measure predictive performance, empirical selection error, intervention sensitivity and computational cost. Together, they connect a formal tolerance to evidence withdrawal with the prediction and maintenance behavior of neural knowledge systems.
